\documentclass[a4paper,11pt]{article}
\usepackage{graphicx}

\DeclareGraphicsExtensions{.pdf,.png,.jpg,.gif}

\newcommand{\HRule}{\rule{\linewidth}{0.5mm}}

%define the title page

\begin{document}

\begin{titlepage}
\begin{center}

\includegraphics[width=0.15\textwidth]{./nsitlogo}~\\[1cm]

\textsc{\LARGE Netaji Subhas Institute of Technology}\\[1.5cm]
\vfill
% Title
\HRule \\[0.4cm]
{\huge \bfseries Intellectual Property Rights \\[0.4cm] }

\HRule \\[1.5cm]

\vfill

% Submission details
\begin{minipage}[t]{0.4\textwidth}
\begin{flushleft} \large
\emph{Submitted By:} \\
Shashank Singh
\end{flushleft}
\end{minipage}
\begin{minipage}[t]{0.4\textwidth}
\begin{flushright} \large
\emph{Submitted to:} \\
Mr.~A.~V.~ Mulley
\end{flushright}
\end{minipage}

\vfill

% Bottom of the page
{\large \today}

\end{center}
\end{titlepage}

\newpage
\tableofcontents

\newpage

\section{Basic Definition}
Intellectual property (IP) refer to creations of the mind, such as inventions,
literary and artistic works, designs and symbols, names and images used in
commerce.

IP is protected in law by, for example \emph{patents}\cite{patents},
\emph{copyright}\cite{copyright} and \emph{trademarks}\cite{trademarks} which
enable people to earn recognition or financial benifit from what they invent or
create. By striking the right balance between the interests of innovators and
wider public interest, the IP system aims to foster an environment in which
creativity and innovation can flourish.

\section{History}
Although many of the legal principles governing IP and IPR have evolved over
centuries, it was not until the 19\textsuperscript{th} century that the term \emph{intellectual
property} began to be used, and not until the late 20\textsuperscript{th}
century that it became commonplace in the majority of the world.\cite{ipweb}
\emph{The Statue of Monopolies (1624)}\cite{statueofmonopolies} and the British
\emph{Statue of Anne (1710)}\cite{statueofanne} are now seen as the origin of
\emph{patent law}\cite{patentlaw} and \emph{copyright}\cite{copyrightlaw}
respectively,\cite{ipweb} firmly establishing the concept of intellectual
property.

The history of patents does not begin with inventions, but rather with royal
grants by \emph{Queen Elizabeth I (1558 - 1603}\cite{queenelizabethi} for monopoly privileges.
Approximately 200 years after the end of Elizabeth's reign, however a
patent represents a legal \emph{right}\cite{legalright} obtained by an inventor
providing for exclusive control over the production and sale of his mechanical
or scientific invention demonstrating the evolution of patents from royal
prerogative to common - law doctrine.

\section{Branches of Intellectual Property}
Intellectual Property is usually divided into two branches, namely
\emph{industrial property}\cite{industrialpropertylaw} and
\emph{copyright}\cite{copyrightlaw}.

\subsection{Industrial Property}
Industrial property typically consists of patents to protect inventions and
industrial designs, which are aesthetic creations determining the appearence of
industrial products. It also covers \emph{trademarks}\cite{trademarks}, trade names and desgnations, as
well as geographical indications, and protection against unfair competition.
They are described as follows:

\subsubsection{Trademarks}
A trademark is a sign capable of distinguishing the goods or services of one
enterprise from those of other enterprise. Trademarks are protected by
\emph{intellectual property}\cite{ipweb} rights.

\subsubsection{Trade names}
Trade names are used by profit and non - profit  entities, political and
religious organizations, industry and agriculture, manufacturers and producers,
wholesalers and retailers, partnerships and coorporations and a host of other
business association. A trade name may be the actual name of a given business
or an assumed name under which a business operates and holds itself out to the
pulic.\cite{tradename}

\subsubsection{Geographical indications}
A geographical indication (GI) is a name used on certain products which
corresponds to a specific geographical location or origin (e.g. a town, region,
or country). The use of GI may act as a certification that the product
possesses certain qualities, is made up according to traditional methods, or
enjoys a certain reputation, due to its geographical
origin.\cite{geographicalindications}

\subsection{Copyright}
Copyright relates to artistic creations, such as poems, novels, music,
paintings, and cinematographic works. In most European languages other than
English, copyright is known as author's right. The expression copyright refers to the main
act which, in respect of literary and artistic creations, may be made only by the
author or with his authorization. That act is the making of copies of the literary
or artistic work, such as a book, a painting, a sculpture, a photograph, or a
motion picture. The second expression, author’s rights refers to the person who 
is the creator of the artistic work, its author, thus underlining the fact, recognized
in most laws, that the author has certain specific rights in his creation, such as 
the right to prevent a distorted reproduction, which only he can exercise, whereas
other rights, such as the right to make copies, can be exercised by other persons,
for example, a publisher who has obtained a license to this effect from the author.

\clearpage
\newpage 

\begin{thebibliography}{12}

\bibitem{patents}
\textit{http://www.wipo.int/patents/en/}

\bibitem{tradename}
\textit{http://legal-dictionary.thefreedictionary.com/trade+name}

\bibitem{geographicalindications}
\textit{http://encyclopedia.thefreedictionary.com/Geographical+Indications}

\bibitem{copyright}
\textit{http://www.wipo.int/copyright/en/}

\bibitem{trademarks}
\textit{http://www.wipo.int/trademarks/en/}

\bibitem{ipweb}
\textit{http://en.wikipedia.org/wiki/Intellectual\textunderscore property}

\bibitem{statueofmonopolies}
\textit{http://en.wikipedia.org/wiki/Statute\textunderscore of\textunderscore Monopolies}

\bibitem{statueofanne}
\textit{http://en.wikipedia.org/wiki/Statute\textunderscore of\textunderscore Anne}

\bibitem{patentlaw}
\textit{http://en.wikipedia.org/wiki/Patent}

\bibitem{copyrightlaw}
\textit{http://en.wikipedia.org/wiki/Copyright}

\bibitem{queenelizabethi}
\textit{http://en.wikipedia.org/wiki/Elizabeth\textunderscore I\textunderscore of\textunderscore England}

\bibitem{legalright}
\textit{http://en.wikipedia.org/wiki/Rights}

\end{thebibliography}
\end{document}
